\section{Reflection}\label{part:reflection}

The preceding sections described CP-SAT from several angles: bounds, propagation,
conflict analysis, lazy encoding, LP reasoning, cuts, restarts, and the parallel
portfolio. The unifying architectural statement is:

\begin{quote}
\textbf{CP-SAT is a CDCL SAT solver whose variables include integer bounds and
whose learnable deductions are generated on demand by explainable propagators,
LP reasoning, and cuts.}
\end{quote}

This statement separates the main ingredients. From constraint programming,
CP-SAT takes \emph{propagation}: specialized algorithms that prune domains using
the structure of constraints such as scheduling, routing, and resource
constraints. From SAT, it takes \emph{learning, backjumping, and restarts}: every
failure can become a reusable clause, and search can jump directly to the
decision level where that clause becomes useful. Lazy clause generation is the
interface that lets these mechanisms meet: propagators produce bound deductions
together with reasons, and those reasons are turned into clauses only when
conflict analysis needs them.

% Sources (the developers' own framing): Laurent Perron, "CP-SAT for scheduling"
% -- "CDCL allows investing more time in costly techniques that benefit only on a
% subset of problems"; "Clause Learning is critical as it offsets the cost of
% running the simplex"; and "[the engine] could incorporate different relaxation /
% specialized (expensive) propagation algorithms, provided they can be EXPLAINED,
% or used at the root node."
Two design principles explain why the hybrid is effective.

% Source (developers' own cost breakdown): Laurent Perron, CPAIOR 2020 Master
% Class (perron2020cpaior), ~17:54 in the main talk (NOT the scheduling-talk
% Q&A) -- "my intuition maybe ... is that we spent one third of the time in the
% [SAT] solver one third of [the] time in the simplex and the rest in
% propagation ... it completely depends on the type of [the] problem but on the
% [SAT] solver it's mostly the conflict analysis which is expensive ... on
% average i would say one third pure sat engine one for cp one third lp." This
% is the empirical anchor for "learning amortizes expensive reasoning": the
% solver deliberately budgets a third of its time on the LP.
\textbf{First, learning amortizes expensive reasoning.} An LP solve, a cutting
plane, or an energetic scheduling deduction is far more expensive than unit
propagation, and on many branches it produces no immediate payoff. It becomes
worth running when the information it produces can be reused. CDCL provides that
reuse mechanism: a costly deduction, once reduced to a learned clause or to an
exact bound reason, can affect many later branches, including branches far from
the one where the deduction was discovered. Thus the relevant question is not
only whether a technique is cheap locally, but whether the information it
produces repays its cost over the whole search tree.

This explains why CP-SAT can place a simplex and a cut loop inside search
workers. In a classical CP solver, such machinery would often be too expensive
to justify at many nodes. In CP-SAT, the information can be banked as bounds,
conflicts, and learned clauses. The developers' rule-of-thumb cost split makes
this bargain concrete: on a typical run, roughly one third of the time is spent
in the SAT engine, one third in the simplex, and one third in propagation. The
costliest part of the SAT third is conflict analysis itself, which is precisely
the mechanism that turns failures into reusable information.

\textbf{Second, explainability is the price of integration.} Any deduction that
changes the in-search trail must be explainable in a common language: Boolean
literals and integer-bound literals. Native propagators provide such reasons
directly. LP reasoning and cuts obey the same rule through exact integer
certificates, as described in \cref{sec:lp,sec:cuts}. The
floating-point LP may suggest a deduction, but only the certified integer reason
is allowed to enter the trail and later conflict analysis.

A technique that cannot produce such a reason is not useless, but it occupies a
different role. It may guide branching, strengthen the model at the root, or
generate candidate solutions in an incomplete worker. It may not directly enqueue
a learnable in-search consequence. This constraint is what lets heterogeneous
techniques behave as parts of one solver rather than as loosely coupled
subsystems.

The resulting integration has several degrees. Native propagators and the
integer trail are tightly coupled: they share decision levels, backtracking, and
conflict analysis. The LP and cut machinery are coupled through certified
reasons: they propose global numeric deductions, then translate the valid ones
into the same reason language. The portfolio is looser still: workers run
separate searches over separate trails, but exchange learned clauses, incumbents,
and bounds. Thus ``one solver'' does not mean one monolithic search tree; it
means one learning currency.

This is the main lesson of CP-SAT's design. It succeeds on many scheduling,
routing, packing, and configuration problems because those domains offer both
rich structure for propagation and deep search spaces in which failures recur.
Propagation exploits the structure; learning prevents the solver from paying for
the same failure repeatedly; LP reasoning and cuts add global numeric bounds; and
the portfolio hedges across strategies. The strength is not any single component
in isolation, but the fact that all components can feed information into the
same learning system.
